\section{Strict vanishing for quartic characters inside the actual two-power host}
\label{sec:cubic-conductor16-strict-vanishing}

\paragraph{Status and conventions.}
The class input in this argument is the separately certified
AC-G142-class.  Its provisional PARI class group is not an input.
All local conditions and coefficient rings below are original ones.
Put $p=269$, $K=\mathbb Q(\omega)$ and $t=\sqrt[3]{p}$ with
its real embedding.  Fix $\chi:G_K\longrightarrow\mu_3$ for
$K(t)/K$, and let $\phi$ be the arithmetic quartic character
modulo sixteen with
\[
 \ker\phi=\{1,7\},\qquad \phi(13)=i,\quad
 \phi(5)=-i,\quad i\equiv82\pmod p.
 \tag{C16S.1}
\]
It is odd and its square is $\chi_8$.  Write $c_p$ for the
actual cyclotomic character and $\omega_p$ for its Teichmuller
part.  The strict group means
\[
 \operatorname{Sel}_{\rm str}(K,D)=\ker\left(
 H^1_{\rm cts}(G_K,D)\longrightarrow
 \prod_{v\mid p}H^1(G_{K_v},D)\times
 \prod_{v\nmid p,\infty}H^1(I_v,D)\right).
 \tag{C16S.2}
\]
In particular the condition at $p$ is zero in full local
cohomology.

\begin{theorem}[AC-G142: conductor-sixteen original strict vanishing]
\label{AC-G142}
Let $E/\mathbb Q_p$ be finite, let $O$ be its original ring of
integers, and assume it contains the values of $\chi$.
For $j=1,2$, $\sigma=\pm1$, $r\in\mathbb Z$, $e\geq1$, and
any continuous $\kappa:G_K\to1+\mathfrak m_O$ unramified away
from $p$, set
\[
 D=(E/O)(\chi^{-j}\phi^{-\sigma}\omega_p^{-r}
                       \kappa^{-1}c_p^{-e}).
 \tag{C16S.3}
\]
If $e+r\equiv267\pmod{268}$, then
\[
                  \operatorname{Sel}_{\rm str}(K,D)=0.
 \tag{C16S.4}
\]
The corresponding residual strict group is zero, over every
finite residue field containing the character values.  No
absolute ramification restriction on $O$ is imposed.
In particular (C16S.4) applies to the four characters
$\chi^j\phi^\sigma$ occurring in the actual AC-K25 host,
at every weight $e\equiv267\pmod{268}$.
\end{theorem}

\begin{proof}
\proofstep{1. The exact quartic field and local decomposition.}
Let $z=\zeta_{16}$ and
\[
 w=z-z^{-1},\quad F=\mathbb Q(w),\quad L=F(t),\quad M=KL.
 \tag{C16S.5}
\]
Then $w^2=-2+\sqrt2$ and $w^4+4w^2+2=0$.
The subgroup $\{1,7\}$ fixes $w$, and its four distinct
conjugates show that $F$ is exactly the field of (C16S.1).
The generator corresponding to $13$ acts by
\[
              b(w)=-w^3-3w,\qquad b^2(w)=-w.
 \tag{C16S.6}
\]
Its unique quadratic subfield is $\mathbb Q(\sqrt2)$.
The $S_3$ splitting field of $x^3-p$ has unique abelian
subfield $K$, which ramifies at three.  Since $F$ is
unramified there, the two Galois fields are disjoint, so
\[
 \operatorname{Gal}(M/\mathbb Q)=S_3\times C_4,
 \qquad\Delta=\operatorname{Gal}(M/K)=C_3\times C_4.
 \tag{C16S.7}
\]
In particular $\mu_p\not\subset M$.

The polynomial $w^4+4w^2+2$ is irreducible modulo $269$;
equivalently the Frobenius class $269\equiv13\pmod{16}$
generates $\operatorname{Gal}(F/\mathbb Q)$.
Thus $F$ is unramified of degree four at $p$, and its
completion contains the unramified quadratic field $K_p$.
Adjoining $t$ is Eisenstein and gives ramification degree
three.  There are two places of $M$ above $p$, each with
absolute ramification degree three and residue degree four.
The decomposition subgroup in $\Delta$ is
$C_3\times\langle b^2\rangle$; the element $b$ exchanges
the two places.  We always use the direct sum over both
completions, not a purported $K_p$-linear action of $b$
on one of them.

\proofstep{2. Inverse-character Kummer descent.}
Let $k=O/\mathfrak m_O$.  The residual character in (C16S.3)
is
\[
 a=\bar\chi^{-j}\bar\phi^{-\sigma}\bar c_p^{-(e+r)},
 \qquad a/\bar c_p=\bar\chi^{-j}\bar\phi^{-\sigma}
                         =:\eta.
 \tag{C16S.8}
\]
The last equality uses $e+r+1\equiv0\pmod{268}$.
The character $\eta$ is faithful on $\Delta$ and has exact
splitting field $M$.  Define
\[
 \begin{split}
 V_{M,S}&=\{u\in M^\times:p\mid v(u)\text{ for }v\nmid p\}/M^{\times p},\\
 V_M&=\{u\in M^\times:p\mid v(u)\text{ for every finite }v\}/M^{\times p},\\
 \mathcal L_M&=\bigoplus_{v\mid p}M_v^\times/M_v^{\times p}.
 \end{split}\tag{C16S.9}
\]
Prime-to-$p$ inflation--restriction and Kummer theory,
with their local counterparts, identify the residual strict
group with
\[
 e_{\eta^{-1}}\ker\left(V_{M,S}\otimes_{\mathbb F_p}k
                  \longrightarrow\mathcal L_M\otimes_{\mathbb F_p}k\right).
 \tag{C16S.10}
\]
The inverse character is forced because the coefficient
restricts to $k(1)$ on $G_M$ and its Kummer cohomology has
the additional twist $\eta$.  Away from $p$, the unramified
condition is divisibility of the Kummer valuation by $p$.
At $p$, a local $p$th power has valuation divisible by $p$,
so the kernel in (C16S.10) belongs to $V_M$.
Its exact ordinary-unit sequence is
\[
 0\longrightarrow\mathcal O_M^\times/\mathcal O_M^{\times p}
 \longrightarrow V_M\longrightarrow\operatorname{Cl}(M)[p]
 \longrightarrow0.
 \tag{C16S.11}
\]

AC-G142-class gives $\operatorname{Cl}(L)[p]=0$.
For the reflection $s$ fixing $L$ in $M$, extension and norm
show that $\operatorname{Cl}(M)[p]^s=0$: their composites
are multiplication by two and $1+s$.  Since $p\ne2$,
$s$ acts as $-1$ on the whole $p$-torsion class group.
Every conjugate reflection does likewise.  Products of
two distinct reflections generate $C_3$, which therefore
acts trivially.  The nontrivial cubic component of the
right term in (C16S.11) is zero.

\proofstep{3. The unit sectors to be tested are simple.}
The twelve complex places of $M$ over the complex place of
$K$ form a regular $\Delta$-set.  Dirichlet's equivariant
unit theorem gives the rational unit character
$\mathbb Q[\Delta]-\mathbb Q$.
Since $p\nmid12$ and $\mu_p(M)=1$, its reduction modulo
$p$ has the same character multiplicities.  For each
quartic eigenvalue $\lambda=82,187$, its nontrivial cubic
part $W_\lambda$ is two dimensional over $\mathbb F_p$.
It is simple as an $\mathbb F_p[C_3]$-module because
$p\equiv2\pmod3$.  Thus a nonzero value of the equivariant
localization on $W_\lambda$ proves injectivity there.
After scalar extension to $k$, this proves injectivity
on both cubic eigenlines.  A fundamental unit basis is
not required for this step.

\proofstep{4. One exact global unit has both nonzero local images.}
The full rational coefficient vector of a unit $U\in L$
and its monic degree-twelve relation are provided in
\path{evidence/G142_sources/root/unit4_exact.json}.
The coefficient basis is $t^a w^c$, $0\leq a\leq2$,
$0\leq c\leq3$, ordered by $4a+c$.
Its SHA256, split for printing, is
\begin{quote}\ttfamily\small
1461afdc3b56ee1fa73404cba0ded8b02\\
d98970c5874f57590b47b3dcbad233f.
\end{quote}
This explicit vector, not a unit-basis assertion from
\texttt{bnfinit}, specifies $U$.  Exact multiplication
matrices with $t^3=p$ and $w^4=-4w^2-2$ give its printed
monic integral characteristic polynomial with constant
coefficient one.  Cayley--Hamilton proves integrality,
and the reversed relation proves that $U^{-1}$ is integral.
Thus $U$ is an actual global unit.  All its tensor
coefficients have denominators prime to $p$.

At one chosen $p$-place use $t$ as uniformizer and identify
the residue field with
$\mathbb F_p[w]/(w^4+4w^2+2)$.  Direct reduction of that
exact vector gives
\[
 U=A+tB\pmod{t^2},\quad
 A=26+232w^2,\qquad B=143+123w+169w^2+39w^3.
 \tag{C16S.12}
\]
Here $A\ne0$.  Let $c$ generate the cubic subgroup and put
$V=U^3/(U\,c(U)\,c^2(U))$.  Since $c(t)=\omega t$,
\[
 V=1+tR\pmod{t^2},\qquad
 R=3B/A=263+135w+266w^2+252w^3.
 \tag{C16S.13}
\]
For each $\lambda=82,187$ form the global unit product
$\prod_{a=0}^3 b^a(V)^{n_a}$, where the integers $n_a$
represent $\lambda^{-a}$ modulo $p$.
Its linear coefficient at the chosen place is
\[
 \begin{array}{c|c}
 \lambda&\sum_{a=0}^3\lambda^{-a}R(b^a w)\\\hline
 82&151w+73w^3\\
 187&120w+128w^3.
 \end{array}\tag{C16S.14}
\]
These are nonzero in the degree-four residue field.
All operations in (C16S.12)--(C16S.14) are in the specified
finite polynomial quotient; the exact script and receipts
are saved with the coefficient vector.

The product is taken globally before evaluation at this
place.  Substituting $b^a w$ therefore does not presume
that $b$ acts on a single completion.  Modulo $p$th powers
the operator $3-(1+c+c^2)$ selects the nontrivial cubic
part, and $\sum_a\lambda^{-a}b^a$ selects the quartic
eigenvalue $\lambda$; the omitted factors three and four
are invertible.  A local $p$th power unit has no nonzero
$t$ coefficient modulo $t^2$: writing it as a residue
lift times $1+tq$, the binomial linear term is divisible
by $pt=t^4$, and all remaining terms have valuation at
least two.  The unramified coefficient field provides
the required residue lifts.
Thus (C16S.14) proves nonzero localization on both
$W_\lambda$.  Step3 proves injectivity on every character
$\bar\chi^j\bar\phi^\sigma$.  Equations
(C16S.10)--(C16S.11) prove residual strict vanishing.

\proofstep{5. Lifting on the original coefficient DVR.}
Let $\varpi$ be a uniformizer of the given $O$.
At the place above $p$ the residual cubic inertia has
order three.  The residual cyclotomic factor has order
dividing $268$, $\phi$ is unramified there, and
$\bar\kappa=1$.  These factors cannot cancel the cubic
inertia.  Hence $D^{G_{K_p}}=0$, and also $D^{G_K}=0$.
Away from $p$, $\kappa$ and $c_p$ are unramified and all
remaining inertia characters have finite order prime
to $p$.  Their invariant module is either zero (a
nontrivial residual inertia value differs from one by
a unit) or all of $D$, which is $\varpi$-divisible.
Thus at every relevant local subgroup $J$,
\[
                       D^J/\varpi D^J=0.
 \tag{C16S.15}
\]
Apply the long exact sequence of
$0\to D[\varpi]\to D\xrightarrow{\varpi}D\to0$
globally and locally.  Any strict class killed by
$\varpi$ lifts globally to $H^1(K,D[\varpi])$;
(C16S.15) makes that lift strict at each place.
The residual strict group is zero, so no nonzero strict
class is killed by $\varpi$.
Every continuous cocycle from the compact group $G_K$
to the discrete module $D$ has finite image, hence is
killed by some power of $\varpi$.  A nonzero strict
class would consequently yield a nonzero one killed
by $\varpi$, a contradiction.  This proves (C16S.4)
on the original $O$, including ramified $O$.
\end{proof}

\paragraph{Actual-host and geometric scope.}
The original growing field is
\[
 H_x=\mathbb Q(\zeta_{2^{h+2}},\zeta_7,
 \sqrt{-q},\sqrt{-d},\omega,t).
\]
For $h\geq2$, it contains $M$, so the four characters
in the theorem occur among the actual AC-K25 characters.
The transfer to the all-prime index gate, including its
local comparison costs, is AC-G143; it is not inferred
by setting individual determinant indices to zero.
No mixed theta realization or cubic RH is asserted here.

In the geometric language this is a statement about
rank-one Artin--Tate coefficient sheaves on the arithmetic
curve of $K$, finite covers split by $M/K$, their ordinary
Kummer unit/class sequence, and restriction to all local
stalk cohomology groups.  The finite-cover and Kummer
dictionaries are the retained SGA~1, Expos\'e~V,
\S\S4 and 7 cover-to-representation dictionary in AC-K25,
and the Stacks Project Kummer sequence, tag~03PK, used in
AC-G103, AC-G114 and AC-G130.  Their earlier primary-source
check scopes are retained; no fresh whole-source proof
inspection is claimed here.
No new Grothendieck archive attribution or purity theorem
is used.  The only new arithmetic inputs are the separately
proved class certificate and the explicit unit calculation
above.
